\section*{Step 315 Attempted Worst-Case-\(\alpha\) Lower Bound}

\paragraph{Requested theorem form.}
The requested argument would infer
\[
|\delta_{D,k}(\rho_1,G_\star)|
\ge |T_{k,j_\ast}|\exp(-\pi^2 k/8),
\]
using \(|\alpha|\le \pi\) and \(\sigma^2\approx k/4\).

\paragraph{Blocker.}
This implication is not rigorous.  A worst-case bound
\(|\alpha|\le\pi\) does not give a positive lower bound on the magnitude
of a complex sum.  Even terms with phases in an interval of length
\(\le 2\pi\) can cancel exactly.  A Gaussian-Fourier heuristic gives
\(\exp(-\alpha^2\sigma^2/2)\) only after one has proved a Gaussian
amplitude profile, an approximately linear phase, and a controlled
stationary-phase remainder.

There is also a direction issue: from \(\sigma^2\ge k/4\) one obtains
\[
\exp(-\pi^2\sigma^2/2)\le \exp(-\pi^2 k/8),
\]
not a lower bound by \(\exp(-\pi^2 k/8)\).  A lower bound of that shape
would require an upper bound \(\sigma^2\le k/4\), or a separate direct
interference estimate.

\paragraph{Finite-sample diagnostic.}
The formal expression
\[
|T_{k,j_\ast}|\exp(-\pi^2k/8)
\]
is below the certified \(|\delta_{D,k}|\) for \(k=10,20,30,50\).  This
is a numerical diagnostic only; it is not a proof for all \(k\).

\paragraph{Conclusion.}
No fully rigorous constants \(A_0,b_0,k_0\) are derived in this step.
The Step 310 theorem remains conditional.  The tightness gap remains the
same: proving a uniform interference lower bound requires ratio-
conjecture-grade or stationary-phase-grade phase control, not only
\(|\alpha|\le\pi\).
